\section{Tổng và hiệu của hai véc-tơ}

\subsection{Đề 1}
\Opensolutionfile{ans}[ans/ans-0-C4B2-1]
\caulc
\begin{ex}%[0H5N2-2]
	Cho $M$ là trung điểm của đoạn thẳng $AB$. Khẳng định nào sau đây là khẳng định đúng?
	\choice
	{$\overrightarrow{IA}+\overrightarrow{IB}=\overrightarrow{AB}$ với $I$ là điểm bất kì}
	{\True $\overrightarrow{AM}+\overrightarrow{BM}=\overrightarrow{0}$}
	{$\overrightarrow{IA}+\overrightarrow{IB}=\overrightarrow{IM}$ với $I$ là điểm bất kì}
	{$\overrightarrow{AM}+\overrightarrow{MB}=\overrightarrow{0}$}
	\loigiai{
		Vì $M$ là trung điểm của $AB$ nên ta luôn có $\overrightarrow{AM}+\overrightarrow{BM}=\overrightarrow{0}$.}
\end{ex}

\begin{ex}%[0H5H2-2]
	Cho hình bình hành $ABCD$ có tâm $O$. Khẳng định nào sau đây \textbf{sai}?
	\choice
	{\True $\overrightarrow{OB}+\overrightarrow{OD}=\overrightarrow{BD}$}
	{$\overrightarrow{OA}+\overrightarrow{OC}=\overrightarrow{0}$}
	{$\overrightarrow{AB}=\overrightarrow{DC}$}
	{$\overrightarrow{AB}+\overrightarrow{AD}=\overrightarrow{AC}$}
	\loigiai{
		\immini{
			Ta có $O$ là trung điểm của $BD\Rightarrow\overrightarrow{OB}+\overrightarrow{OD}=\overrightarrow{0}$.\\
			Do đó khẳng định $\overrightarrow{OB}+\overrightarrow{OD}=\overrightarrow{BD}$ là sai.
		}{\begin{tikzpicture}[line join=round,line cap=round,>=stealth,scale=.75,font=\footnotesize]
				\foreach \x/\y/\n in {0/0/A,1/2/B,4/0/D} \coordinate (\n) at (\x,\y);
				\coordinate (C) at ($(B)+(D)-(A)$); 
				\coordinate (O) at ($(B)!0.5!(D)$); 
				\draw (A)--(B)--(C)--(D)--(A)--(C) (B)--(D);
				\foreach \t/\g in {A/-135,C/45,B/135,D/-45,O/-90}{
					\draw[fill=black] (\t) circle (1pt) node[shift={(\g:8pt)}]{$ \t $};
				}
		\end{tikzpicture}}
	}
\end{ex}

\begin{ex}%[0H5H2-2]
	Đẳng thức nào sau đây luôn đúng với mọi điểm $A$, $B$, $C$ bất kì?
	\choice
	{$\overrightarrow{AB}+\overrightarrow{CB}=\overrightarrow{AC}$}
	{\True $\overrightarrow{CB}+\overrightarrow{AC}=\overrightarrow{AB}$}
	{$\overrightarrow{AB}-\overrightarrow{AC}=\overrightarrow{BC}$}
	{$\overrightarrow{BC}-\overrightarrow{AC}=\overrightarrow{AB}$}
	\loigiai{
		Ta có $\overrightarrow{CB}+\overrightarrow{AC}=\overrightarrow{AC}+\overrightarrow{CB}=\overrightarrow{AB}$.}
\end{ex}

\begin{ex}%[0H5H2-2]
	Cho ba điểm $A$, $B$, $C$ phân biệt. Đẳng thức nào sau đây là đẳng thức \textbf{sai}?
	\choice
	{$\overrightarrow{BA}+\overrightarrow{AC}=\overrightarrow{BC}$}
	{$\overrightarrow{AB}+\overrightarrow{BC}=\overrightarrow{AC}$}
	{\True $\overrightarrow{CA}+\overrightarrow{AB}=\overrightarrow{BC}$}
	{$\overrightarrow{AB}-\overrightarrow{AC}=\overrightarrow{CB}$}
	\loigiai{
		Theo quy tắc ba điểm, ta có  $\overrightarrow{CA}+\overrightarrow{AB}=\overrightarrow{CB}$ nên $\overrightarrow{CA}+\overrightarrow{AB}=\overrightarrow{BC}$ là sai.}
\end{ex}

\begin{ex}%[0H5H2-2]
	Gọi $O$ là tâm hình bình hành $ABCD$. Đẳng thức nào sau đây \textbf{sai}?
	\choice
	{$\overrightarrow{AB}-\overrightarrow{AD}=\overrightarrow{DB}$}
	{$\overrightarrow{BC}-\overrightarrow{BA}=\overrightarrow{DC}-\overrightarrow{DA}$}
	{$\overrightarrow{OA}-\overrightarrow{OB}=\overrightarrow{CD}$}
	{\True $\overrightarrow{OB}-\overrightarrow{OC}=\overrightarrow{OD}-\overrightarrow{OA}$}
	\loigiai{
		Ta có 
		$\overrightarrow{OB}-\overrightarrow{OC}=\overrightarrow{OD}-\overrightarrow{OA}\Leftrightarrow\overrightarrow{CB}=\overrightarrow{AD}$.\\
		Đáp án này sai vì $\overrightarrow{CB}=\overrightarrow{DA}$ (do $ABCD$ là hình bình hành).}
\end{ex}

\begin{ex}%[0H5H2-2]
	Chọn khẳng định \textbf{sai}?
	\choice
	{Nếu $I$ là trung điểm đoạn $AB$ thì $\overrightarrow{AI}+\overrightarrow{IB}=\overrightarrow{AB}$}
	{\True Nếu $I$ là trung điểm đoạn $AB$ thì $\overrightarrow{IA}+\overrightarrow{BI}=\overrightarrow{0}$}
	{Nếu $I$ là trung điểm đoạn $AB$ thì $\overrightarrow{AI}+\overrightarrow{BI}=\overrightarrow{0}$}
	{Nếu $I$ là trung điểm đoạn $AB$ thì $\overrightarrow{IA}+\overrightarrow{IB}=\overrightarrow{0}$}
	\loigiai{
		Ta có $\overrightarrow{IA}+\overrightarrow{BI}=\overrightarrow{BI}+\overrightarrow{IA}=\overrightarrow{BA}$.\\
		Vậy $\overrightarrow{IA}+\overrightarrow{BI}=\overrightarrow{0}$ là sai.}
\end{ex}

\begin{ex}%[0H5N2-2]
	Cho tam giác $ABC$. Khẳng định nào sau đây là \textbf{sai}?
	\choice
	{$\overrightarrow{AB}-\overrightarrow{AC}=\overrightarrow{CB}$}
	{$\overrightarrow{AB}+\overrightarrow{BA}=\overrightarrow{0}$}
	{\True $\overrightarrow{AB}-\overrightarrow{AC}=\overrightarrow{BC}$}
	{$\overrightarrow{AB}+\overrightarrow{BC}=\overrightarrow{AC}$}
	\loigiai{
		Ta có $\overrightarrow{AB}-\overrightarrow{AC}=\overrightarrow{CB}$ nên khẳng định $\overrightarrow{AB}-\overrightarrow{AC}=\overrightarrow{BC}$ là sai.}
\end{ex}

\begin{ex}%[0H5H2-2]
	Gọi $G$ là trọng tâm của tam giác $ABC$ và $M$ là điểm bất kỳ. Chọn khẳng định \textbf{sai} trong các khẳng định sau.
	\choice
	{$\overrightarrow{AG}+\overrightarrow{BG}+\overrightarrow{CG}=\overrightarrow{0}$}
	{$\overrightarrow{MA}+\overrightarrow{MB}+\overrightarrow{MC}=3\overrightarrow{MG}$}
	{$\overrightarrow{GA}+\overrightarrow{GB}+\overrightarrow{GC}=\overrightarrow{0}$}
	{\True $\overrightarrow{MA}+\overrightarrow{MB}+\overrightarrow{MC}=\overrightarrow{MG}$}
	\loigiai{
		Với $G$ là trọng tâm của tam giác $ABC$ và $M$ là điểm bất kỳ, ta có
		\begin{itemize}
			\item $\overrightarrow{AG}+\overrightarrow{BG}+\overrightarrow{CG}=\overrightarrow{0}\Leftrightarrow-\overrightarrow{GA}-\overrightarrow{GB}-\overrightarrow{GC}=\overrightarrow{0}\Leftrightarrow\overrightarrow{GA}+\overrightarrow{GB}+\overrightarrow{GC}=\overrightarrow{0}$.
			\item  $\overrightarrow{GA}+\overrightarrow{GB}+\overrightarrow{GC}=\overrightarrow{0}$.
			\item $\overrightarrow{MA}+\overrightarrow{MB}+\overrightarrow{MC}=3\overrightarrow{MG}$.
		\end{itemize}
	}
\end{ex}

\begin{ex}%[0H5H2-1]
	Cho hình bình hành $ABCD$ tâm $O$. Khi đó $\overrightarrow{OA}+\overrightarrow{BO}$ bằng
	\choice
	{$\overrightarrow{OC}+\overrightarrow{OB}$}
	{$\overrightarrow{AB}$}
	{\True $\overrightarrow{CD}$}
	{$\overrightarrow{OC}+\overrightarrow{DO}$}
	\loigiai{
		\immini{
			Ta có $\overrightarrow{OA}+\overrightarrow{BO}=\overrightarrow{BO}+\overrightarrow{OA}=\overrightarrow{BA}$.\\
			Vì $ABCD$ là hình bình hành nên $\overrightarrow{BA}=\overrightarrow{CD}$.
		}{\begin{tikzpicture}[line join=round,line cap=round,>=stealth,scale=.75,font=\footnotesize]
				\foreach \x/\y/\n in {0/0/A,1/2/B,4/0/D} \coordinate (\n) at (\x,\y);
				\coordinate (C) at ($(B)+(D)-(A)$); 
				\coordinate (O) at ($(B)!0.5!(D)$); 
				\draw (A)--(B)--(C)--(D)--(A)--(C) (B)--(D);
				\foreach \t/\g in {A/-135,C/45,B/135,D/-45,O/-90}{
					\draw[fill=black] (\t) circle (1pt) node[shift={(\g:8pt)}]{$ \t $};
				}
		\end{tikzpicture}}	
	}
\end{ex}

\begin{ex}%[0H5H2-2]
	Cho hình vuông $ABCD$, có $M$ là giao điểm của hai đường chéo. Trong các mệnh đề sau, tìm mệnh đề \textbf{sai}?
	\choice
	{$\overrightarrow{AB}+\overrightarrow{BC}=\overrightarrow{AC}$}
	{$\overrightarrow{AB}+\overrightarrow{AD}=\overrightarrow{AC}$}
	{$\overrightarrow{BA}+\overrightarrow{BC}=2\overrightarrow{BM}$}
	{\True $\overrightarrow{MA}+\overrightarrow{MB}=\overrightarrow{MC}+\overrightarrow{MD}$}
	\loigiai{
		\immini{
			Ta có 
			$$\begin{aligned}
				&\overrightarrow{MA}+\overrightarrow{MB}=\overrightarrow{MC}+\overrightarrow{MD}\\
				\Leftrightarrow\ & \overrightarrow{MA}-\overrightarrow{MC}=\overrightarrow{MD}-\overrightarrow{MB}\\
				\Leftrightarrow\ & \overrightarrow{CA}=\overrightarrow{BD}.
			\end{aligned}$$
			Vì $\overrightarrow{CA}=\overrightarrow{BD}$ sai nên $\overrightarrow{MA}+\overrightarrow{MB}=\overrightarrow{MC}+\overrightarrow{MD}$ là sai.
		}{\begin{tikzpicture}[line join=round,line cap=round,>=stealth,scale=.75,font=\footnotesize]
				\foreach \x/\y/\n in {0/0/A,0/3/B,3/0/D} \coordinate (\n) at (\x,\y);
				\coordinate (C) at ($(B)+(D)-(A)$); 
				\coordinate (M) at ($(B)!0.5!(D)$); 
				\draw (A)--(B)--(C)--(D)--(A)--(C) (B)--(D);
				\foreach \t/\g in {A/-135,C/45,B/135,D/-45,M/-90}{
					\draw[fill=black] (\t) circle (1pt) node[shift={(\g:8pt)}]{$ \t $};
				}
		\end{tikzpicture}}
	}
\end{ex}

\begin{ex}%[0H5H2-2]
	Cho hình bình hành $ABCD$ tâm $O$. Đẳng thức nào sau đây đúng?
	\choice
	{\True $\overrightarrow{OA}+\overrightarrow{OC}=\overrightarrow{0}$}
	{$\overrightarrow{OA}+\overrightarrow{CO}=\overrightarrow{0}$}
	{$\overrightarrow{AO}+\overrightarrow{OC}=\overrightarrow{0}$}
	{$\overrightarrow{OA}-\overrightarrow{OC}=\overrightarrow{0}$}
	\loigiai{
		\immini{
			Vì $ABCD$ là hình bình hành nên $O$ là trung điểm của $AC$.\\
			Vậy $\overrightarrow{OA}+\overrightarrow{OC}=\overrightarrow{0}$.
		}{\begin{tikzpicture}[line join=round,line cap=round,>=stealth,scale=.75,font=\footnotesize]
				\foreach \x/\y/\n in {0/0/A,1/2/B,4/0/D} \coordinate (\n) at (\x,\y);
				\coordinate (C) at ($(B)+(D)-(A)$); 
				\coordinate (O) at ($(B)!0.5!(D)$); 
				\draw (A)--(B)--(C)--(D)--(A)--(C) (B)--(D);
				\foreach \t/\g in {A/-135,C/45,B/135,D/-45,O/-90}{
					\draw[fill=black] (\t) circle (1pt) node[shift={(\g:8pt)}]{$ \t $};
				}
		\end{tikzpicture}}	
	}
\end{ex}

\begin{ex}%[0H5H2-2]
	Cho lục giác đều $ABCDEF$ và $O$ là tâm của nó. Đẳng thức nào dưới đây là đẳng thức \textbf{sai}?
	\choice
	{$\overrightarrow{OA}+\overrightarrow{OC}+\overrightarrow{OE}=\overrightarrow{0}$}
	{$\overrightarrow{BC}+\overrightarrow{FE}=\overrightarrow{AD}$}
	{$\overrightarrow{OA}+\overrightarrow{OB}+\overrightarrow{OC}=\overrightarrow{EB}$}
	{\True $\overrightarrow{AB}+\overrightarrow{CD}+\overrightarrow{FE}=\overrightarrow{0}$}
	\loigiai{
		\immini{
			Ta có $\overrightarrow{AB}+\overrightarrow{CD}+\overrightarrow{FE}=\overrightarrow{AB}+\overrightarrow{BO}+\overrightarrow{AO}=2\overrightarrow{AO}\ne\overrightarrow{0}$.\\
			Do đó mệnh đề $\overrightarrow{AB}+\overrightarrow{CD}+\overrightarrow{FE}=\overrightarrow{0}$ là sai.
		}{\begin{tikzpicture}[line join=round,line cap=round,>=stealth,scale=.75,font=\footnotesize]
				\coordinate (O) at (0,0);
				\foreach \n/\p in {A/120,B/60,C/0,D/-60,E/-120,F/-180} \coordinate (\n) at ($(O)+({2*cos(\p)},{2*sin(\p)})$);
				\draw (A)--(B)--(C)--(D)--(E)--(F)--(A)--(D) (B)--(E) (C)--(F);
				\foreach \t/\g in {A/120,B/60,C/0,D/-60,E/-120,F/-180,O/-90}{
					\draw[fill=black] (\t) circle (1pt) node[shift={(\g:8pt)}]{$ \t $};
				}
		\end{tikzpicture}}		
	}
\end{ex}

\Closesolutionfile{ans}
\indapan{6}{ans/ans-0-C4B2-1}
\cauds
\Opensolutionfile{ans}[ans/ans-0-C4B2-1-DS]

\begin{ex}%[0H5V2-2]
	Cho hình bình hành $ABCD$. Xét tính đúng sai của các mệnh đề sau.
	\choiceTF
	{\True $\overrightarrow{C B}+\overrightarrow{C D}=\overrightarrow{C A}$}
	{$\overrightarrow{AC}+\overrightarrow{DA}=\overrightarrow{CD}$}
	{$\overrightarrow{BA}-\overrightarrow{BC}=\overrightarrow{AC}$}
	{\True $\overrightarrow{BA}-\overrightarrow{BC}+\overrightarrow{AD}=\overrightarrow{CD}$}
	\loigiai{
		\begin{itemchoice}
			\itemch Đúng. Vì 
			$\overrightarrow{CB}+\overrightarrow{CD}=\overrightarrow{CA}$ (quy tắc hình bình hành).
			\itemch Sai. Vì  $\overrightarrow{AC}+\overrightarrow{DA}=\overrightarrow{DC}$ (quy tắc ba điểm).
			\itemch Sai. Vì $\overrightarrow{BA}-\overrightarrow{BC}=\overrightarrow{CA}$ (quy tắc ba điểm).
			\itemch Đúng. Vì 
			$$\begin{aligned}
				\overrightarrow{BA}-\overrightarrow{BC}+\overrightarrow{AD}&=\overrightarrow{BA}+\overrightarrow{CB}+\overrightarrow{AD}\\
				&=\left(\overrightarrow{CB}+\overrightarrow{BA}\right)+\overrightarrow{AD}\\
				&=\overrightarrow{CA}+\overrightarrow{AD}=\overrightarrow{CD} \text{ (quy tắc ba điểm)}.
			\end{aligned}$$ 
		\end{itemchoice}
	}
\end{ex}

\begin{ex}%[0H5V2-2]
	Cho hình bình hành $ABCD$ tâm $O$. Xét tính đúng sai của các mệnh đề sau.
	\choiceTF
	{\True $\overrightarrow{BA}+\overrightarrow{AB}=\vec{0}$}
	{\True $\overrightarrow{OA}+\overrightarrow{AC}=\overrightarrow{OC}$}
	{$\overrightarrow{OA}-\overrightarrow{OB}+\overrightarrow{DC}=\overrightarrow{CB}$}
	{\True $\overrightarrow{OB}+\overrightarrow{OA}-\overrightarrow{CA}=\overrightarrow{DC}$}
	\loigiai{
		\begin{itemchoice}
			\itemch Đúng. Vì $\overrightarrow{BA}+\overrightarrow{AB}=\overrightarrow{BB}=\vec{0}$.
			\itemch Đúng. Vì $\overrightarrow{OA}+\overrightarrow{AC}=\overrightarrow{OC}$ (quy tắc ba điểm).
			\itemch Sai. Vì $\overrightarrow{OA}-\overrightarrow{OB}+\overrightarrow{DC}=\overrightarrow{OA}+\overrightarrow{BO}+\overrightarrow{DC}=\overrightarrow{BO}+\overrightarrow{OA}+\overrightarrow{AB}=\overrightarrow{BA}+\overrightarrow{AB}=\overrightarrow{0}$.
			\itemch Đúng. Vì $\overrightarrow{OB}+\overrightarrow{OA}-\overrightarrow{CA}=\overrightarrow{OB}+\overrightarrow{OA}+\overrightarrow{AC}=\overrightarrow{OB}+\overrightarrow{OC}=\overrightarrow{DO}+\overrightarrow{OC}=\overrightarrow{DC}$.
		\end{itemchoice}
	}
\end{ex}

\begin{ex}%[0H5V2-2]
	Cho sáu điểm $A$, $B$, $C$, $D$, $E$, $F$. Xét tính đúng sai của các mệnh đề sau.
	\choiceTF
	{$\overrightarrow{AB}+\overrightarrow{CD}+\overrightarrow{EF}-\overrightarrow{CB}-\overrightarrow{ED}=\overrightarrow{FA}$}
	{$\overrightarrow{AB}-\overrightarrow{AF}+\overrightarrow{CD}-\overrightarrow{CB}+\overrightarrow{EF}=\overrightarrow{DE}$}
	{\True $\overrightarrow{AB}+\overrightarrow{CD}=\overrightarrow{AD}+\overrightarrow{CB}$}
	{\True $\overrightarrow{AC}+\overrightarrow{BD}+\overrightarrow{EF}=\overrightarrow{A F}+\overrightarrow{BC}+\overrightarrow{ED}$}
	\loigiai{
		\begin{itemchoice}
			\itemch Sai. Vì 
			$$\begin{aligned}
				\overrightarrow{AB}+\overrightarrow{CD}+\overrightarrow{EF}-\overrightarrow{CB}-\overrightarrow{ED}
				&=\overrightarrow{AB}+\overrightarrow{CD}+\overrightarrow{EF}+\overrightarrow{BC}+\overrightarrow{DE}\\
				&=\left(\overrightarrow{AB}+\overrightarrow{BC}\right)+\left(\overrightarrow{CD}+\overrightarrow{DE}\right)+\overrightarrow{EF}\\
				&=\overrightarrow{AC}+\overrightarrow{CE}+\overrightarrow{EF}=\overrightarrow{AE}+\overrightarrow{EF}=\overrightarrow{AF}.
			\end{aligned}$$
			\itemch Sai. Vì $\overrightarrow{AB}-\overrightarrow{AF}+\overrightarrow{CD}-\overrightarrow{CB}+\overrightarrow{EF}=\overrightarrow{FB}+\overrightarrow{BD}+\overrightarrow{EF}=\overrightarrow{FD}+\overrightarrow{EF}=\overrightarrow{ED}$.
			\itemch Đúng. Vì $\overrightarrow{AB}+\overrightarrow{CD}=\overrightarrow{AD}+\overrightarrow{CB}\Leftrightarrow\overrightarrow{AB}-\overrightarrow{AD}=\overrightarrow{CB}-\overrightarrow{CD}\Leftrightarrow\overrightarrow{DB}=\overrightarrow{DB}$.
			\itemch Đúng. Vì 
			$$\begin{aligned}
				&\overrightarrow{AC}+\overrightarrow{BD}+\overrightarrow{EF}=\overrightarrow{AF}+\overrightarrow{BC}+\overrightarrow{ED}\\
				\Leftrightarrow\ & \left(\overrightarrow{AC}-\overrightarrow{AF}\right)+\left(\overrightarrow{BD}-\overrightarrow{BC}\right)+\left(\overrightarrow{EF}-\overrightarrow{ED}\right)=\vec{0}\\ 
				\Leftrightarrow\ &\overrightarrow{FC}+\overrightarrow{CD}+\overrightarrow{DF}=\vec{0}\Leftrightarrow\overrightarrow{FD}+\overrightarrow{DF}=\vec{0}. 
			\end{aligned}$$
		\end{itemchoice}
	}
\end{ex}

\begin{ex}%[0H5V2-4]
	Cho hình vuông $ABCD$ có cạnh bằng $a$. Xét tính đúng sai của các mệnh đề sau.
	\choiceTF
	{\True $\overrightarrow{AD}+\overrightarrow{AB}=\overrightarrow{AC}$}
	{\True Gọi $E$ là điểm đối xứng với $B$ qua $C$. Khi đó $ADEC$ là hình thang}
	{\True $\left|\overrightarrow{AD}+\overrightarrow{AB}\right|=a\sqrt{2}$}
	{$\left|\overrightarrow{AD}+\overrightarrow{AC}\right|=a\sqrt{3}$}
	\loigiai{
		\begin{itemchoice}
			\itemch Đúng. Vì theo quy tắc hình bình hành ta có $\overrightarrow{AD}+\overrightarrow{AB}=\overrightarrow{AC}$.
			\itemch Đúng. Vì với $E$ là điểm đối xứng với $B$ qua $C$, ta có $CE=AD$ và $CE\parallel AD$.
			\immini{Suy ra $ADEC$ là hình bình hành, hay $ADEC$ cũng là hình thang.\\
			}{\begin{tikzpicture}[line join=round,line cap=round,>=stealth,scale=.7,font=\footnotesize]
					\foreach \x/\y/\n in {0/0/A,0/3/B,3/0/D} \coordinate (\n) at (\x,\y);
					\coordinate (C) at ($(B)+(D)-(A)$); 
					\coordinate (E) at ($(C)+(C)-(B)$); 
					\draw (A)--(B)--(C)--(D)--(A)--(C) (B)--(D)--(E)--(C);
					\foreach \t/\g in {A/-135,C/45,B/135,D/-45,E/45}{
						\draw[fill=black] (\t) circle (1pt) node[shift={(\g:8pt)}]{$ \t $};
					}
			\end{tikzpicture}}
			\itemch Đúng. Vì theo định lí Pytago ta có $AC=\sqrt{AB^2+BC^2}=\sqrt{a^2+a^2}=a\sqrt{2}$.\\
			Vậy $\left|\overrightarrow{AD}+\overrightarrow{AB}\right|=\left|\overrightarrow{AC}\right|=AC=a\sqrt{2}$.
			\itemch Sai. Vì ta có $\overrightarrow{AD}+\overrightarrow{AC}=\overrightarrow{AE}$. Áp dụng định lí Pytago cho tam giác $ABE$, ta có\\
			$AE^2=AB^2+BE^2=a^2+\left(2a\right)^2=5a^2\Rightarrow AE=a\sqrt{5}$.\\
			Vậy $\left|\overrightarrow{AD}+\overrightarrow{AC}\right|=AE=a\sqrt{5}$.
		\end{itemchoice}
	}
\end{ex}

\Closesolutionfile{ans}
\indapan{3}{ans/ans-0-C4B2-1-DS}

\Opensolutionfile{ans}[ans/ans-0-C4B2-1-KQ]
\caukq

\begin{ex}%[0H5H2-4]
	Cho hình vuông $ABCD$ có độ dài cạnh bằng $2$, $M$ là trung điểm $BC$. Tính $\left|\overrightarrow{AB}+\overrightarrow{BM}\right|$ (kết quả làm tròn đến hàng phần trăm).
	\shortans{$2{,}24$}
	\loigiai{
		\immini{
			Ta có $\left|\overrightarrow{AB}+\overrightarrow{BM}\right|=\left|\overrightarrow{AM}\right|=AM$.\\
			Theo định lí Py-ta-go, ta có $ AM=\sqrt{AB^2+BM^2}=\sqrt{2^2+1^2}=\sqrt{5}$.\\
			Vậy $\left|\overrightarrow{AB}+\overrightarrow{BM}\right|=AM=\sqrt{5}\approx 2{,}24$.
		}{\begin{tikzpicture}[line join=round,line cap=round,>=stealth,scale=.7,font=\footnotesize]
				\foreach \x/\y/\n in {0/0/A,0/3/B,3/0/D} \coordinate (\n) at (\x,\y);
				\coordinate (C) at ($(B)+(D)-(A)$); 
				\coordinate (M) at ($(C)!0.5!(B)$); 
				\draw (A)--(B)--(C)--(D)--(A)--(M);
				\foreach \t/\g in {A/-135,C/45,B/135,D/-45,M/90}{
					\draw[fill=black] (\t) circle (1pt) node[shift={(\g:8pt)}]{$ \t $};
				}
		\end{tikzpicture}}	
	}
\end{ex}

\begin{ex}%[0H5V2-6]
	Cho hai lực $\vec{F}_1$, $\vec{F}_2$ có điểm đặt $O$ tạo với nhau góc $60^{\circ}$, biết rằng cường độ của hai lực $\vec{F}_1$ và $\vec{F}_2$ đều bằng $100\sqrt{3}$ N. Tính cường độ tổng hợp của hai lực trên?
	\shortans{$200$}
	\loigiai{
		\immini{
			Chọn các điểm $A$, $B$ thỏa mãn $\vec{F}_1=\overrightarrow{OA}$, $\vec{F}_2=\overrightarrow{OB}$.\\ 
			Gọi điểm $C$ là một đỉnh của hình bình hành $OACB$, khi đó \begin{center}
				$\vec{F}_1+\vec{F}_2=\overrightarrow{OA}+\overrightarrow{OB}=\overrightarrow{OC}$ (quy tắc hình bình hành).
			\end{center}
			Cường độ tổng hợp hai lực là $\left|\vec{F}_1+\vec{F}_2\right|=\left|\overrightarrow{OC}\right|=OC$.
		}{
			\begin{tikzpicture}[line join=round,line cap=round,>=stealth,scale=1,font=\footnotesize]
				\foreach \x/\y/\n in {0/0/O} \coordinate (\n) at (\x,\y);
				\coordinate (A) at ($(O)+({2*cos(30)},{2*sin(30)})$);
				\coordinate (B) at ($(O)+({2*cos(-30)},{2*sin(-30)})$);
				\coordinate (C) at ($(A)+(B)-(O)$);
				\coordinate (I) at ($(A)!0.5!(B)$);
				\draw[arrow style,,thick] (O)--(A) node[sloped,midway,above]{$\vec{F}_1$};
				\draw[arrow style,,thick] (O)--(B) node[sloped,midway,below]{$\vec{F}_2$};
				\draw[arrow style,dashed,thick] (O)--(C);
				\draw[dashed] (A)--(C)--(B)--(A);
				\foreach \t/\g in {A/30,B/-30,O/180,C/0,I/-45}{
					\draw[fill=black] (\t) circle (1pt) node[shift={(\g:10pt)}]{$ \t $};
				}
				\draw pic[thin,draw,"$60^\circ$",angle radius=.8cm] {angle= B--O--A};
				\draw pic[thin,draw,angle radius=2mm] {right angle= A--I--O};
		\end{tikzpicture}}
		\noindent Xét tam giác $OAB$ có $OA=OB=100\sqrt{3}$ và $\widehat{AOB}=60^\circ$ nên tam giác $OAB$ đều.\\
		Gọi $I$ là tâm hình bình hành $OACB$, khi đó $OI$ cũng là đường cao tam giác đều $OAB$.\\
		Do đó $OI=\dfrac{100\sqrt{3}\cdot\sqrt{3}}{2}=100$, suy ra $OC=2OI=200$.\\
		Vậy hợp lực của $\vec{F}_1$, $\vec{F}_2$ có độ lớn là $200$ N.
	}
\end{ex}

\begin{ex}%[0H5V2-4]
	Cho hình vuông $ABCD$ có độ dài cạnh bằng $2\sqrt{2}$. Tính $\left|\overrightarrow{AB}+\overrightarrow{AC}+\overrightarrow{AD}\right|$.
	\shortans{$8$}
	\loigiai{
		\immini{
			Ta có $\overrightarrow{AB}+\overrightarrow{AC}+\overrightarrow{AD}=(\overrightarrow{AB}+\overrightarrow{AD})+\overrightarrow{AC}=\overrightarrow{AC}+\overrightarrow{AC}$.\\
			Gọi $E$ đối xứng với $A$ qua $C$, suy ra $\overrightarrow{AC}=\overrightarrow{CE}$.\\
			Khi đó $\overrightarrow{AB}+\overrightarrow{AC}+\overrightarrow{AD}=\overrightarrow{AC}+\overrightarrow{AC}=\overrightarrow{AC}+\overrightarrow{CE}=\overrightarrow{AE}$.\\
			Ta có $AE=2AC=2\cdot 2\sqrt{2}\cdot\sqrt{2}=8$.\\
			Vậy $\left|\overrightarrow{AB}+\overrightarrow{AC}+\overrightarrow{AD}\right|=AE=8$.
		}{\begin{tikzpicture}[line join=round,line cap=round,>=stealth,scale=.6,font=\footnotesize]
				\foreach \x/\y/\n in {0/0/A,0/2/B,2/0/D} \coordinate (\n) at (\x,\y);
				\coordinate (C) at ($(B)+(D)-(A)$); 
				\coordinate (E) at ($(A)!2!(C)$); 
				\draw (A)--(B)--(C)--(D)--(A)--(E);
				\foreach \t/\g in {A/-135,C/0,B/135,D/-45,E/45}{
					\draw[fill=black] (\t) circle (1pt) node[shift={(\g:8pt)}]{$ \t $};
				}
		\end{tikzpicture}}	
	}
\end{ex}

\begin{ex}%[0H5H2-4]
	Cho tam giác vuông $ABC$ có các cạnh góc vuông là $AB=1$, $AC=2$.
	Điểm $M$ thỏa mãn $\overrightarrow{AC}-\overrightarrow{AB}=\overrightarrow{AM}$. Tính độ dài véc-tơ $\overrightarrow{AM}$ (kết quả là tròn đến hàng phần chục).
	\shortans{$2{,}2$}
	\loigiai{
		\immini{
			Ta có $\overrightarrow{AC}-\overrightarrow{AB}=\overrightarrow{AM}\Leftrightarrow \overrightarrow{BC}=\overrightarrow{AM}$.\\
			Suy ra $ABCM$ là hình bình hành.\\
			Ta có $\left|\overrightarrow{AM}\right|=AM=BC=\sqrt{1^2+2^2}=\sqrt{5}\approx 2{,}2$.
		}{\begin{tikzpicture}[line join=round,line cap=round,>=stealth,scale=1,font=\footnotesize]
				\foreach \x/\y/\n in {0/0/A,2/0/C,0/1/B} \coordinate (\n) at (\x,\y);
				\coordinate (M) at ($(A)+(C)-(B)$);  
				\draw (A)--(B)--(C)--(M)--(A)--(C);
				\foreach \t/\g in {A/-135,C/45,M/-45,B/135}{
					\draw[fill=black] (\t) circle (1pt) node[shift={(\g:8pt)}]{$ \t $};
				}
				\draw pic[thin,draw,angle radius=2mm] {right angle= B--A--C};
		\end{tikzpicture}}
	}
\end{ex}

\begin{ex}%[0H5V2-6]
	Cho hai lực $\vec{F}_1$, $\vec{F}_2$ có điểm đặt $A$ tạo với nhau góc $45^{\circ}$, biết rằng cường độ của hai lực $\vec{F}_1$ và $\vec{F}_2$ lần lượt bằng $60$ N, $90$ N. Tính cường độ tổng hợp của hai lực trên (đáp án làm tròn đến hàng đơn vị của Newton).
	\shortans{$139$}
	\loigiai{
		\immini{
			Đặt $\vec{F}_1=\overrightarrow{AB}$, $\vec{F}_2=\overrightarrow{AD}$.\\
			Vẽ hình bình hành $ABCD$.\\
			Ta có $\vec{F}_1+\vec{F}_2=\overrightarrow{AB}+\overrightarrow{AD}=\overrightarrow{AC}$.\\
			Vì $\widehat{BAD}=45^{\circ}\Rightarrow\widehat{ABC}=135^{\circ}$; $AD=90=BC$.\\
			Theo định lí cosin ta có
		}{\begin{tikzpicture}[line join=round,line cap=round,>=stealth,scale=1.1,font=\footnotesize]
				\foreach \x/\y/\n in {0/0/A} \coordinate (\n) at (\x,\y);
				\coordinate (B) at ($(A)+({2*cos(45)},{2*sin(45)})$);
				\coordinate (D) at ($(A)+({3*cos(0)},{3*sin(0)})$);
				\coordinate (C) at ($(B)+(D)-(A)$);
				\draw[arrow style,thick] (A)--(B) node[sloped,midway,above]{$\vec{F}_1$};
				\draw[arrow style,thick] (A)--(D) node[sloped,midway,below]{$\vec{F}_2$};
				\draw[arrow style,dashed,thick] (A)--(C);
				\draw[dashed] (B)--(C)--(D);
				\foreach \t/\g in {B/45,D/-45,A/-135,C/30}{
					\draw[fill=black] (\t) circle (1pt) node[shift={(\g:10pt)}]{$ \t $};
				}
				\draw pic[thin,draw,"",angle radius=.5cm] {angle= D--A--B};
				\draw pic[thin,draw,"$135^\circ$",angle radius=.6cm] {angle= A--B--C};
		\end{tikzpicture}}
		\noindent
		$$AC=\sqrt{AB^2+BC^2-2AB\cdot BC\cdot\cos135^{^\circ}}=\sqrt{60^2+90^2-2\cdot 60\cdot 90\cdot\left(-\dfrac{\sqrt{2}}{2}\right)}\approx 139.$$
		Vậy véc-tơ hợp lực của $\vec{F}_1$, $\vec{F}_2$ có độ lớn là $\left|\vec{F}_1+\vec{F}_2\right|\approx 139$ N.
	}
\end{ex}

\begin{ex}%[0H5V2-6]
	Một dòng sông chảy từ phía Bắc xuống phía Nam với vận tốc $10$ km/h, có một chiếc ca nô chuyển động từ phía Đông sang phía Tây với vận tốc $35$ km/h so với dòng nước. Vận tốc của ca nô so với bờ là bao nhiêu km/h (kết quả làm tròn đến hàng phần chục)?
	\shortans{$36{,}4$}
	\loigiai{
		\immini{
			Gọi $\vec{v}_1$, $\vec{v}_2$ lần lượt là vectơ vận tốc của dòng nước đối với bờ và ca nô đối với dòng nước. Khi đó vận tốc của ca nô đối với bờ chính là tổng $\vec{v}_1+\vec{v}_2$.\\ 
			Đặt $\vec{v}_1=\overrightarrow{AD}$, $\vec{v}_2=\overrightarrow{AB}$ với $A$ là vị trí của ca nô.\\
			Vẽ hình bình hành $ABCD$, ta có $$\vec{v}_1+\vec{v}_2=\overrightarrow{AB}+\overrightarrow{AD}=\overrightarrow{AC}.$$
			Theo định lí Py-ta-go $AC=\sqrt{10^2+35^2}=5\sqrt{53}\approx 36{,}4$.\\
			Vậy vận tốc của ca nô đối với bờ là xấp xỉ $36{,}4$ km/h.
		}{
			\begin{tikzpicture}[line join=round,line cap=round,>=stealth,scale=1.1,font=\footnotesize]
				\foreach \x/\y/\n in {0/0/4} \coordinate (\n) at (\x,\y);
				\coordinate (B) at ($(A)+({1*cos(-90)},{1*sin(-90)})$);
				\coordinate (D) at ($(A)+({3.5*cos(180)},{3.5*sin(180)})$);
				\coordinate (C) at ($(B)+(D)-(A)$);
				\coordinate (M) at ($(D)!0.5!(A)$);
				\coordinate (o) at ($(M)+(0,2)$);
				\coordinate (b) at ($(o)+(0,.5)$);
				\coordinate (t) at ($(o)+(-.5,0)$);
				\coordinate (d) at ($(o)+(.5,0)$);
				\coordinate (n) at ($(o)+(0,-.5)$);
				\draw[arrow style,thick] (A)--(B) node[sloped,midway,above]{$\vec{v}_1$};
				\draw[arrow style,thick] (A)--(D) node[sloped,midway,above]{$\vec{v}_2$};
				\draw[arrow style,dashed,thick] (A)--(C);
				\draw[dashed] (B)--(C)--(D);
				\draw[arrow style] (o)--(b)node[above]{Bắc};
				\draw[arrow style] (o)--(d)node[right]{Đông};
				\draw[arrow style] (o)--(n)node[below]{Nam};
				\draw[arrow style] (o)--(t)node[left]{Tây};
				\foreach \t/\g in {B/-45,D/135,A/45,C/-135}{
					\draw[fill=black] (\t) circle (1pt) node[shift={(\g:10pt)}]{$ \t $};
				}
				\draw pic[thin,draw,angle radius=2mm] {right angle= B--A--D};
		\end{tikzpicture}}
	}
\end{ex}

\Closesolutionfile{ans}
\indapan{6}{ans/ans-0-C4B2-1-KQ}
\newpage 
\subsection{Đề 2}
\Opensolutionfile{ans}[ans/ans-0-C4B2-2]
\caulc
\begin{ex}%[0H5H2-2]
	Gọi $O$ là tâm hình bình hành $ABCD$; hai điểm $E$, $F$ lần lượt là trung điểm $AB$, $BC$. Đẳng thức nào sau đây \textbf{sai}?
	\choice
	{$\overrightarrow{DO}=\overrightarrow{EB}-\overrightarrow{EO}$}
	{$\overrightarrow{OC}=\overrightarrow{EB}+\overrightarrow{EO}$}
	{$\overrightarrow{OA}+\overrightarrow{OC}+\overrightarrow{OD}+\overrightarrow{OE}+\overrightarrow{OF}=\vec{0}$}
	{\True $\overrightarrow{BE}+\overrightarrow{BF}-\overrightarrow{DO}=\vec{0}$}
	\loigiai{
	\immini{
		Ta có $OF$, $OE$ lần lượt là đường trung bình của tam giác $BCD$ và $ABC$.\\
		Suy ra $BEOF$ là hình bình hành.\\
		Khi đó $\overrightarrow{BE}+\overrightarrow{BF}=\overrightarrow{BO}$.\\
		Suy ra $\overrightarrow{BE}+\overrightarrow{BF}-\overrightarrow{DO}=\overrightarrow{BO}-\overrightarrow{DO}=\overrightarrow{OD}-\overrightarrow{OB}=\overrightarrow{BD}$.
		}{\begin{tikzpicture}[line join=round,line cap=round,>=stealth,scale=1,font=\footnotesize]
			\foreach \x/\y/\n in {0/0/A,3/0/D,-1/-2/B} \coordinate (\n) at (\x,\y);
			\coordinate (C) at ($(B)+(D)-(A)$);
			\coordinate (E) at ($(A)!1/2!(B)$);
			\coordinate (F) at ($(B)!1/2!(C)$); 
			\coordinate (O) at ($(A)!1/2!(C)$);     
			\draw (A)--(B)--(C)--(D)--(A)--(C) (B)--(D) (O)--(E)--(F)--(O);
			\foreach \t/\g in {A/90,B/-135,C/-45,D/45,E/180,F/-90,O/90}{
				\draw[fill=black] (\t) circle (1pt) node[shift={(\g:8pt)}]{$ \t $};
			}
			\end{tikzpicture}}
		}
\end{ex}

\begin{ex}%[0H5H2-2]
	Cho hình bình hành $ABCD$. Gọi $G$ là trọng tâm của tam giác $ABC$. Mệnh đề nào sau đây đúng?
	\choice
	{\True $\overrightarrow{GA}+\overrightarrow{GC}+\overrightarrow{GD}=\overrightarrow{BD}$}
	{$\overrightarrow{GA}+\overrightarrow{GC}+\overrightarrow{GD}=\overrightarrow{CD}$}
	{$\overrightarrow{GA}+\overrightarrow{GC}+\overrightarrow{GD}=\overrightarrow{O}$}
	{$\overrightarrow{GA}+\overrightarrow{GD}+\overrightarrow{GC}=\overrightarrow{CD}$}
	\loigiai{
		\immini{
		Vì $G$ là trọng tâm của $\triangle ABC$ nên $\overrightarrow{GA}+\overrightarrow{GB}+\overrightarrow{GC}=\overrightarrow{0}$.\\
		Do đó 
		$$\begin{aligned}
			\overrightarrow{GA}+\overrightarrow{GC}+\overrightarrow{GD}&=\overrightarrow{GA}+\overrightarrow{GC}+\left(\overrightarrow{GB}+\overrightarrow{BC}+\overrightarrow{CD}\right)\\
			&=\left(\overrightarrow{GA}+\overrightarrow{GB}+\overrightarrow{GC}\right)+\overrightarrow{BC}+\overrightarrow{CD}\\
			&=\overrightarrow{BC}+\overrightarrow{CD}=\overrightarrow{BD}.
		\end{aligned}$$
	}{\begin{tikzpicture}[line join=round,line cap=round,>=stealth,scale=1,font=\footnotesize]
		\foreach \x/\y/\n in {0/0/A,3/0/D,-1/-2/B} \coordinate (\n) at (\x,\y);
		\coordinate (C) at ($(B)+(D)-(A)$);
		\coordinate (M) at ($(B)!1/2!(C)$); 
		\coordinate (O) at ($(A)!1/2!(C)$);     
		\coordinate (G) at ($(B)!2/3!(O)$);  
		\draw (A)--(B)--(C)--(D)--(A)--(C) (B)--(D) (A)--(M);
		\foreach \t/\g in {A/90,B/-135,C/-45,D/45,G/-30,M/-90,O/90}{
			\draw[fill=black] (\t) circle (1pt) node[shift={(\g:8pt)}]{$ \t $};
		}
		\end{tikzpicture}}
	}
\end{ex}

\begin{ex}%[0H5N2-4]
	Cho tam giác đều $ABC$ có cạnh bằng $a$. Khi đó $\left|\overrightarrow{CB}-\overrightarrow{CA}\right|$ bằng
	\choice
	{$a\sqrt{3}$}
	{$a$}
	{$2a$}
	{\True $\dfrac{a\sqrt{3}}{2}$}
	\loigiai{
		Ta có $\left|\overrightarrow{CB}-\overrightarrow{CA}\right|=\left|\overrightarrow{AB}\right|=a$ do tam giác đều $ABC$ có cạnh bằng $a$.}
\end{ex}

\begin{ex}%[0H5N2-4]
	Cho tam giác $ABC$ đều cạnh $a$. Độ dài $\overrightarrow{AB}+\overrightarrow{BC}$ bằng
	\choice
	{\True $a$}
	{$\dfrac{a\sqrt{3}}{2}$}
	{$2a$}
	{$a\sqrt{3}$}
	\loigiai{
		Ta có 
		$\left|\overrightarrow{AB}+\overrightarrow{BC}\right|=\left|\overrightarrow{AC}\right|=AC=a$.}
\end{ex}

\begin{ex}%[0H5H2-4]
	Cho hình vuông $ABCD$ cạnh $a$, tâm $O$. Tính độ dài của véc-tơ $\overrightarrow{OA}+\overrightarrow{OB}$.
	\choice
	{$2a$}
	{$\dfrac{a}{2}$}
	{\True $a$}
	{$3a$}
	\loigiai{
		\immini{
		Gọi $M$ là trung điểm của $AB$.\\
		Suy ra $\left|\overrightarrow{OA}+\overrightarrow{OB}\right|=2\left|\overrightarrow{OM}\right|=2OM=AD=a$.
	}{\begin{tikzpicture}[line join=round,line cap=round,>=stealth,scale=1,font=\footnotesize]
		\foreach \x/\y/\n in {0/0/A,2/0/B,0/-2/D} \coordinate (\n) at (\x,\y);
		\coordinate (C) at ($(B)+(D)-(A)$);
		\coordinate (O) at ($(B)!0.5!(D)$);
		\coordinate (M) at ($(A)!1/2!(B)$);   
		\draw (A)--(B)--(C)--(D)--(A)--(C) (D)--(B) (M)--(O);
		\foreach \t/\g in {A/135,D/-135,C/-45,B/45,O/-90,M/90}{
			\draw[fill=black] (\t) circle (1pt) node[shift={(\g:8pt)}]{$ \t $};
		}
		\end{tikzpicture}}
		}
\end{ex}

\begin{ex}%[0H5H2-4]
	Cho hình chữ nhật $ABCD$, $AB=3$, $AD=4$. Tính $\left|\overrightarrow{AB}+\overrightarrow{AD}\right|$.
	\choice
	{$\left|\overrightarrow{AB}+\overrightarrow{AD}\right|=8$}
	{$\left|\overrightarrow{AB}+\overrightarrow{AD}\right|=7$}
	{$\left|\overrightarrow{AB}+\overrightarrow{AD}\right|=6$}
	{\True $\left|\overrightarrow{AB}+\overrightarrow{AD}\right|=5$}
	\loigiai{
		\immini{
		Ta có $\overrightarrow{AB}+\overrightarrow{AD}=\overrightarrow{AC}$.\\
		Suy ra $\left|\overrightarrow{AB}+\overrightarrow{AD}\right|=\left|\overrightarrow{AC}\right|=AC=\sqrt{A{B^2}+B{C^2}}=\sqrt{3^2+4^2}=5$.
	}{\begin{tikzpicture}[line join=round,line cap=round,>=stealth,scale=1,font=\footnotesize]
		\foreach \x/\y/\n in {0/0/A,2/0/D,0/1.5/B} \coordinate (\n) at (\x,\y);
		\coordinate (C) at ($(B)+(D)-(A)$);
		\coordinate (O) at ($(B)!0.5!(D)$);
		\coordinate (M) at ($(A)!1/2!(B)$);
		\draw (A)--(B)--(C)--(D)--(A)--(C) (D)--(B);
		\foreach \t/\g in {A/-135,D/-45,C/45,B/135,O/-90}{
			\draw[fill=black] (\t) circle (1pt) node[shift={(\g:8pt)}]{$ \t $};
		}
		\end{tikzpicture}}
		}
\end{ex}

\begin{ex}%[0H5V3-7]
	Cho tam giác $ABC$ cân tại $C$. Tập hợp các điểm $M$ thỏa mãn đẳng thức $\left|\overrightarrow{MA}+\overrightarrow{MB}\right|=2\left|\overrightarrow{MC}\right|$ là
	\choice
	{\True Đường thẳng song song với $AB$}
	{Đường thẳng vuông góc với $ AB$}
	{Một điểm}
	{Một đường tròn}
	\loigiai{
		Gọi $I$ là trung điểm $AB$, ta có
		$$\left|\overrightarrow{MA}+\overrightarrow{MB}\right|=2\left|\overrightarrow{MC}\right|\Leftrightarrow 2\left|\overrightarrow{MI}\right|=2\left|\overrightarrow{MC}\right|\Leftrightarrow MI=MC.$$
		Điều đó chứng tỏ điểm $M$ cách đều hai điểm $I$, $C$, nên tập hợp các điểm $M$ là đường trung trực của đoạn $IC$.\\
		Mặt khác, tam giác $ABC$ cân tại $C$ suy ra $IC\perp AB$ nên đường trung trực của đoạn $IC$ song song với đường thẳng $AB$.}
\end{ex}

\begin{ex}%[0H5H2-4]
	Cho hình vuông $ABCD$ cạnh $a$, tâm $O$. Tính độ dài của véc-tơ $\overrightarrow{OA}+\overrightarrow{OB}$.
	\choice
	{$2a$}
	{$\dfrac{a}{2}$}
	{\True $a$}
	{$3a$}
	\loigiai{
		\immini{Ta có $\overrightarrow{OA}+\overrightarrow{OB}=\overrightarrow{OE}=2\overrightarrow{OI}$,
		với $E$ là đỉnh thứ tư của hình bình hành $AOBE$ và $\{I\}=AB\cap OE$.\\
		Suy ra $\left|\overrightarrow{OA}+\overrightarrow{OB}\right|=\left| 2\overrightarrow{OI}\right|=2\cdot\left|\overrightarrow{OI}\right|=2\cdot OI$.\\
		Trong tam giác vuông $ABD$ ta có 
		\begin{center}
			$OI$ là đường trung bình nên  $OI=\dfrac{AD}{2}=\dfrac{a}{2}$.
		\end{center}
		Vậy $\left|\overrightarrow{OA}+\overrightarrow{OB}\right|=2\cdot\dfrac{a}{2}=a$.
		}{\begin{tikzpicture}[line join=round,line cap=round,>=stealth,scale=1,font=\footnotesize]
		\foreach \x/\y/\n in {0/0/A,2/0/B,0/-2/D} \coordinate (\n) at (\x,\y);
		\coordinate (C) at ($(B)+(D)-(A)$);
		\coordinate (O) at ($(B)!0.5!(D)$);
		\coordinate (E) at ($(A)+(B)-(O)$);
		\coordinate (I) at ($(A)!0.5!(B)$);
		\draw (A)--(B)--(C)--(D)--(A)--(C) (D)--(B)--(E)--(A) (O)--(E);
		\foreach \t/\g in {A/135,D/-135,C/-45,B/45,O/-90,E/90,I/-45}{
			\draw[fill=black] (\t) circle (1pt) node[shift={(\g:8pt)}]{$ \t $};
		}
		\end{tikzpicture}}
		}
\end{ex}

\begin{ex}%[0H5V2-4]
	Cho hình bình hành $ABCD$ có $AB=a$, $AB\perp BD$, $\widehat{BAD}=60^\circ $. Gọi $E$, $F$ lần lượt là trung điểm của $BD$, $AD$. Độ dài véc-tơ $\overrightarrow{BE}+\overrightarrow{AF}$ là
	\choice
	{\True $\dfrac{a\sqrt{13}}{2}$}
	{$\dfrac{a\sqrt{10}}{2}$}
	{$\dfrac{a\sqrt{7}}{2}$}
	{$ 2a$}
	\loigiai{
		\immini{
		Gọi $H$ là trung điểm $AB$.\\
		Ta có $\overrightarrow{BE}=\overrightarrow{ED}$, $\overrightarrow{AF}=\overrightarrow{HE}$ nên suy ra
		$$\overrightarrow{BE}+\overrightarrow{AF}=\overrightarrow{HE}+\overrightarrow{ED}=\overrightarrow{HD}$$
		Suy ra $\left|\overrightarrow{BE}+\overrightarrow{AF}\right|=\left|\overrightarrow{HD}\right|=HD$.\\
		Ta lại có $BD=AB\tan 60^\circ=a\sqrt{3}$.\\ $HD=\sqrt{HB^2+BD^2}=\sqrt{\left(\dfrac{a}{2}\right)^2+\left(a\sqrt{3}\right)^2}=\dfrac{a\sqrt{13}}{2}$.\\
		Vậy $\left|\overrightarrow{BE}+\overrightarrow{AF}\right|=\dfrac{a\sqrt{13}}{2}$.
	}{\begin{tikzpicture}[line join=round,line cap=round,>=stealth,scale=1,font=\footnotesize]
		\foreach \x/\y/\n in {0/0/A,4/0/D} \coordinate (\n) at (\x,\y);
		\coordinate (B) at ($(A)+({2*cos(60)},{2*sin(60)})$);
		\coordinate (C) at ($(B)+(D)-(A)$);
		\coordinate (E) at ($(D)!1/2!(B)$);
		\coordinate (F) at ($(A)!1/2!(D)$); 
		\coordinate (H) at ($(A)!1/2!(B)$);         
		\draw (A)--(B)--(C)--(D)--(A) (B)--(D) (E)--(F) (D)--(H)--(E);
		\foreach \t/\g in {A/225,B/135,C/45,D/-45,E/60,F/-90,H/180}{
			\draw[fill=black] (\t) circle (1pt) node[shift={(\g:8pt)}]{$ \t $};
		}
		\draw pic[thin,draw,angle radius=2mm] {right angle= A--B--D};
		\end{tikzpicture}}
		}
\end{ex}

\begin{ex}%[0H5H2-4]
	Cho hình vuông $ABCD$ cạnh $a$. Tính $\left|\overrightarrow{AB}-\overrightarrow{DA}\right|$ là
	\choice
	{$\left|\overrightarrow{AB}-\overrightarrow{DA}\right|=2a$}
	{$\left|\overrightarrow{AB}-\overrightarrow{DA}\right|=0$}
	{\True $\left|\overrightarrow{AB}-\overrightarrow{DA}\right|=a\sqrt{2}$}
	{$\left|\overrightarrow{AB}-\overrightarrow{DA}\right|=a$}
	\loigiai{
		\immini{
		Ta có $\overrightarrow{AB}=\overrightarrow{DC}$ nên $$\left|\overrightarrow{AB}-\overrightarrow{DA}\right|=\left|\overrightarrow{DC}-\overrightarrow{DA}\right|=\left|\overrightarrow{AC}\right|=AC=a\sqrt{2}.$$
	}{\begin{tikzpicture}[line join=round,line cap=round,>=stealth,scale=1,font=\footnotesize]
		\foreach \x/\y/\n in {0/0/A,2/0/D,0/2/B} \coordinate (\n) at (\x,\y);
		\coordinate (C) at ($(B)+(D)-(A)$);
		\coordinate (O) at ($(B)!0.5!(D)$);
		\coordinate (M) at ($(A)!1/2!(B)$);
		\draw (A)--(B)--(C)--(D)--(A)--(C) (D)--(B);
		\foreach \t/\g in {A/-135,D/-45,C/45,B/135,O/-90}{
			\draw[fill=black] (\t) circle (1pt) node[shift={(\g:8pt)}]{$ \t $};
		}
		\end{tikzpicture}}
		}
\end{ex}

\begin{ex}%[0H5H2-4]
	Cho hình vuông $ABCD$ cạnh $a$. Tính $S=\left| 2\overrightarrow{AD}+\overrightarrow{DB}\right|$?
	\choice
	{$S=a$}
	{$S=a\sqrt{3}$}
	{\True $S=a\sqrt{2}$}
	{$S=a\sqrt{5}$}
	\loigiai{
		\immini{
		Ta có $2\overrightarrow{AD}+\overrightarrow{DB}=\overrightarrow{AD}+\overrightarrow{AB}=\overrightarrow{AC}$.\\
		Suy ra $\left| 2\overrightarrow{AD}+\overrightarrow{DB}\right|=\left|\overrightarrow{AC}\right|=AC=a\sqrt{2}$.\\
		Vậy $S=a\sqrt{2}$.
		}{\begin{tikzpicture}[line join=round,line cap=round,>=stealth,scale=1,font=\footnotesize]
		\foreach \x/\y/\n in {0/0/A,2/0/D,0/2/B} \coordinate (\n) at (\x,\y);
		\coordinate (C) at ($(B)+(D)-(A)$);
		\coordinate (O) at ($(B)!0.5!(D)$);
		\coordinate (M) at ($(A)!1/2!(B)$);
		\draw (A)--(B)--(C)--(D)--(A)--(C) (D)--(B);
		\foreach \t/\g in {A/-135,D/-45,C/45,B/135,O/-90}{
			\draw[fill=black] (\t) circle (1pt) node[shift={(\g:8pt)}]{$ \t $};
		}
		\end{tikzpicture}}
		}
\end{ex}

\begin{ex}%[0H5V2-6]
	Cho ba lực $\vec{F}_1=\overrightarrow{MA}$, $\vec{F}_2=\overrightarrow{MB}$, $\vec{F}_3=\overrightarrow{MC}$ cùng tác động vào một vật tại điểm $M$ và vật đứng yên.
	\begin{center}
		\begin{tikzpicture}[line join=round,line cap=round,>=stealth,scale=1,font=\footnotesize]
			\foreach \x/\y/\n in {0/0/M} \coordinate (\n) at (\x,\y);
			\coordinate (A) at ($(M)+({2*cos(135)},{2*sin(135)})$);
			\coordinate (B) at ($(M)+({2*cos(225)},{2*sin(225)})$);
			\coordinate (D) at ($(A)+(B)-(M)$);
			\coordinate (C) at ($(M)+(M)-(D)$);
			\shade[ball color=white] (M) circle (.2cm);
			\draw[->,thick] (M)--(A) node[sloped,midway,above]{$\vec{F}_1$};
			\draw[->,thick] (M)--(B) node[sloped,midway,below]{$\vec{F}_2$};
			\draw[->,thick] (M)--(C) node[sloped,midway,above]{$\vec{F}_3$};
			\foreach \t/\g in {A/150,B/-150,M/-90,C/0}{
				\draw[fill=black] (\t) circle (1pt) node[shift={(\g:10pt)}]{$ \t $};
			}
			\draw pic[thin,draw,angle radius=3mm] {right angle= A--M--B};
		\end{tikzpicture}
	\end{center}
	Cho biết cường độ của $\vec{F}_1$, $\vec{F}_2$ đều bằng $100$ N và góc $\widehat{AMB}=90^\circ$. Khi đó cường độ của lực $\vec{F}_3$ là
	\choice
	{$50\sqrt{3}$ N}
	{$100\sqrt{3}$ N}
	{$50\sqrt{2}$ N}
	{\True $100\sqrt{2}$ N}
	\loigiai{
		\begin{center}
			\begin{tikzpicture}[line join=round,line cap=round,>=stealth,scale=1,font=\footnotesize]
				\foreach \x/\y/\n in {0/0/M} \coordinate (\n) at (\x,\y);
				\coordinate (A) at ($(M)+({2*cos(135)},{2*sin(135)})$);
				\coordinate (B) at ($(M)+({2*cos(225)},{2*sin(225)})$);
				\coordinate (D) at ($(A)+(B)-(M)$);
				\coordinate (C) at ($(M)+(M)-(D)$);
				\shade[ball color=white] (M) circle (.2cm);
				\draw[->,thick] (M)--(A) node[sloped,midway,above]{$\vec{F}_1$};
				\draw[->,thick] (M)--(B) node[sloped,midway,below]{$\vec{F}_2$};
				\draw[->,thick] (M)--(C) node[sloped,midway,above]{$\vec{F}_3$};
				\draw[->,thick] (M)--(D);
				\draw[dashed] (A)--(D) (B)--(D);
				\foreach \t/\g in {A/150,B/-150,M/-90,C/0,D/180}{
					\draw[fill=black] (\t) circle (1pt) node[shift={(\g:10pt)}]{$ \t $};
				}
				\draw pic[thin,draw,angle radius=3mm] {right angle= A--M--B};
			\end{tikzpicture}
		\end{center}
		Vẽ hình vuông $ADBM$, cạnh $MA=100$ thì $MD=100\sqrt{2}$.\\
		Theo qui tắc hình bình hành có $\overrightarrow{MA}+\overrightarrow{MB}=\overrightarrow{MD}$.\\
		Vật đứng yên khi $\vec{F}_3=-\overrightarrow{MD}$. Suy ra cường độ của lực $\vec{F}_3$ là $\left|\vec{F}_3\right|=100\sqrt{2}N$.}
\end{ex}

\Closesolutionfile{ans}
\indapan{6}{ans/ans-0-C4B2-2}
\cauds
\Opensolutionfile{ans}[ans/ans-0-C4B2-2-DS]

\begin{ex}%[0H5H2-2]
	Cho tam giác $ABC$. Các điểm $M$, $N$, $P$ lần lượt là trung điểm của $AB$, $AC$, $BC$. Xét tính đúng sai của các mệnh đề sau.
	\choiceTF
	{\True $\overrightarrow{A M}-\overrightarrow{A N}=\overrightarrow{NM}$}
	{$\overrightarrow{MN}-\overrightarrow{NC}=\overrightarrow{MP}$}
	{\True $\overrightarrow{MN}-\overrightarrow{PN}=\overrightarrow{MP}$}
	{$\overrightarrow{BP}-\overrightarrow{CP}=\overrightarrow{PC}$}
\loigiai{
	Dễ thấy tứ giác $AMPN$, $MNCP$ là hình bình hành.
	\begin{center}
		\begin{tikzpicture}[line join=round,line cap=round,>=stealth,scale=1,font=\footnotesize]
		\foreach \x/\y/\n in {0/0/B,4/0/C,1/2/A} \coordinate (\n) at (\x,\y);
		\coordinate (M) at ($(A)!1/2!(B)$);
		\coordinate (N) at ($(A)!1/2!(C)$);
		\coordinate (P) at ($(B)!1/2!(C)$);   
		\draw (A)--(B)--(C)--(A) (M)--(N)--(P)--(M);
		\foreach \t/\g in {A/90,B/-135,C/-45,M/180,N/30,P/-90}{
			\draw[fill=black] (\t) circle (1pt) node[shift={(\g:8pt)}]{$ \t $};
		}
	\end{tikzpicture}
	\end{center}
	\begin{itemchoice}
		\itemch Đúng. Vì $\overrightarrow{AM}-\overrightarrow{A N}=\overrightarrow{NM}$.
		\itemch Sai. Vì $\overrightarrow{MN}-\overrightarrow{NC}=\overrightarrow{MN}-\overrightarrow{MP}=\overrightarrow{PN}\ (\overrightarrow{NC}=\overrightarrow{MP})$.
		\itemch Đúng. Vì $\overrightarrow{MN}-\overrightarrow{PN}=\overrightarrow{MN}+\overrightarrow{NP}=\overrightarrow{MP}$.
		\itemch Sai. Vì $\overrightarrow{BP}-\overrightarrow{CP}=\overrightarrow{BP}+\overrightarrow{PC}=\overrightarrow{BC}$.
	\end{itemchoice}
}
\end{ex}

\begin{ex}%[0H5H2-4]
	Cho hình vuông $ABCD$ cạnh $a$, có $O$ là giao điểm hai đường chéo. Xét tính đúng sai của các mệnh đề sau.
	\choiceTF
	{\True $O$ là trung điểm của $AC$, $BD$}
	{$\left|\overrightarrow{OA}-\overrightarrow{CB}\right|=a\sqrt{2}$}
	{$\left|\overrightarrow{AB}+\overrightarrow{DC}\right|=a$}
	{$\left|\overrightarrow{CD}-\overrightarrow{DA}\right|=\dfrac{a\sqrt{2}}{2}$}
\loigiai{
	\begin{itemchoice}
		\itemch Đúng. $O$ là giao điểm hai đường chéo nên $O$ là trung điểm của $AC$, $BD$.
		\itemch Sai. Vì $\left|\overrightarrow{OA}-\overrightarrow{CB}\right|=\left|\overrightarrow{CO}-\overrightarrow{CB}\right|=\left|\overrightarrow{BO}\right|=BO=\dfrac{a\sqrt{2}}{2}$.
		\itemch Sai. Vì $\left|\overrightarrow{AB}+\overrightarrow{DC}\right|=\left|\overrightarrow{AB}+\overrightarrow{AB}\right|=2\left|\overrightarrow{AB}\right|=2a$.
		\itemch Sai. Vì $\left|\overrightarrow{CD}-\overrightarrow{DA}\right|=\left|\overrightarrow{CD}-\overrightarrow{CB}\right|=\left|\overrightarrow{BD}\right|=BD=a\sqrt{2}$.
	\end{itemchoice}
	}
\end{ex}

\begin{ex}%[0H5H2-4]
	Cho hình thoi $ABCD$ cạnh $a$, có $\widehat{BAD}=60^{\circ}$. Gọi $O$ là giao điểm hai đường chéo. Xét tính đúng sai của các mệnh đề sau.
	\choiceTF
	{\True $\triangle BAD$, $\triangle BCD$ đều cạnh $a$}
	{$\left|\overrightarrow{AB}+\overrightarrow{AD}\right|=a\sqrt{2}$}
	{\True $\left|\overrightarrow{BA}-\overrightarrow{BC}\right|=a\sqrt{3}$}
	{\True $\left|\overrightarrow{OB}-\overrightarrow{DC}\right|=\dfrac{a\sqrt{3}}{2}$}
\loigiai{
	\begin{itemchoice}
		\itemch Đúng. Vì giả thiết cho ta $\triangle BAD$, $\triangle BCD$ đều cạnh $a$.
		\itemch Sai. Vì $\left|\overrightarrow{AB}+\overrightarrow{AD}\right|=\left|\overrightarrow{AC}\right|=AC=2AO=2\sqrt{AB^2-BO^2}=a\sqrt{3}$.
		\itemch Đúng. Vì $\left|\overrightarrow{BA}-\overrightarrow{BC}\right|=\left|\overrightarrow{CA}\right|=CA=a\sqrt{3}$.
		\itemch Đúng. Vì $\left|\overrightarrow{OB}-\overrightarrow{DC}\right|=\left|\overrightarrow{DO}-\overrightarrow{DC}\right|=\left|\overrightarrow{CO}\right|=CO=AO=\dfrac{a\sqrt{3}}{2}$.
	\end{itemchoice}
	}
\end{ex}









\begin{ex}%[0H5H2-4]
	Cho tam giác $ABC$ đều có cạnh bằng $3 a$. Xét tính đúng sai của các mệnh đề sau.
	\choiceTF
	{\True $A O=3 a\dfrac{\sqrt{3}}{2}$}
	{$\left|\overrightarrow{AB}-\overrightarrow{AC}\right|=2a$}
	{\True $\left|\overrightarrow{AB}-\overrightarrow{CB}-\overrightarrow{AC}\right|=0$}
	{$\left|\overrightarrow{AB}+\overrightarrow{AC}\right|=a\sqrt{3}$}
\loigiai{
	\begin{itemchoice}
		\itemch Đúng. Vì $AO=3a\dfrac{\sqrt{3}}{2}$.
		\itemch Sai. Vì $\left|\overrightarrow{A B}-\overrightarrow{A C}\right|=|\overrightarrow{CB}|=CB=3a$.
		\itemch Đúng. Vì 
		$\left|\overrightarrow{AB}-\overrightarrow{CB}-\overrightarrow{AC}\right|=\left|\overrightarrow{AB}-\overrightarrow{AC}-\overrightarrow{CB}\right|=\left|\overrightarrow{CB}-\overrightarrow{CB}\right|=\left|\overrightarrow{0}\right|=0$.
		\itemch Sai. Dựng hình bình hành $ABDC$, ta có $\overrightarrow{AB}+\overrightarrow{AC}=\overrightarrow{AD}$.
		\immini{
		Gọi $O$ là giao điểm của $AD$ và $BC$. \\
		Ta có $AO=3 a\dfrac{\sqrt{3}}{2}$.\\
		Suy ra $\left|\overrightarrow{AB}+\overrightarrow{AC}\right|=\left|\overrightarrow{AD}\right|=AD=2AO=3a\sqrt{3}$.
		}{\begin{tikzpicture}[line join=round,line cap=round,>=stealth,scale=1,font=\footnotesize]
			\foreach \x/\y/\n in {0/0/B,2/0/C} \coordinate (\n) at (\x,\y);
			\coordinate (A) at ($(B)+({2*cos(60)},{2*sin(60)})$);
			\coordinate (D) at ($(B)+(C)-(A)$);
			\coordinate (O) at ($(B)!0.5!(C)$);
			\draw (A)--(B)--(C)--(D)--(A)--(C) (B)--(D);
			\foreach \t/\g in {A/90,D/-90,C/0,B/180,O/-45}{
				\draw[fill=black] (\t) circle (1pt) node[shift={(\g:8pt)}]{$ \t $};
			}
			\draw pic[thin,draw,angle radius=2mm] {right angle= A--O--B};
			\end{tikzpicture}}
	\end{itemchoice}
	}
\end{ex}



\Closesolutionfile{ans}
\indapan{3}{ans/ans-0-C4B2-2-DS}

\Opensolutionfile{ans}[ans/ans-0-C4B2-2-KQ]
\caukq

\begin{ex}%[0H5H2-4]
	Cho hình chữ nhật $ABCD$, $AB=3$, $AD=4$. Tính $\left|\overrightarrow{AB}+\overrightarrow{AD}\right|$.
	\shortans{$5$}
\loigiai{
	Ta có $\overrightarrow{AB}+\overrightarrow{AD}=\overrightarrow{AC}$ $\Rightarrow\left|\overrightarrow{AB}+\overrightarrow{AD}\right|=\left|\overrightarrow{AC}\right|=AC=\sqrt{AB^2+BC^2}=5$.}
\end{ex}

\begin{ex}%[0H5H2-6]
	\immini{Cho hai lực $\vec{F}_1=\overrightarrow{MA}$, $\vec{F}_2=\overrightarrow{MB}$ cùng tác động vào một vật tại điểm $M$. Cường độ hai lực $\vec{F}_1$, $\vec{F}_2$ lần lượt là $300$ N và $400$ N, $\widehat{AMB}=90^\circ $. Cường độ của lực tác động lên vật là bao nhiêu?
	}{\begin{tikzpicture}[line join=round,line cap=round,>=stealth,scale=.5,font=\footnotesize]
		\foreach \x/\y/\n in {0/0/M,4/0/A,0/-3/B} \coordinate (\n) at (\x,\y);
		\shade[ball color=white] (M) circle (.3cm);
		\draw[->,thick] (M)--(A);
		\draw[->,thick] (M)--(B);
		\foreach \t/\g in {A/0,B/-90,M/180}{
			\draw[fill=black] (\t) circle (1pt) node[shift={(\g:10pt)}]{$ \t $};
		}
		\draw pic[thin,draw,angle radius=3mm] {right angle= A--M--B};
		\end{tikzpicture}}
	\shortans{$500$}
\loigiai{
	Ta có tổng lực tác dụng lên vật là $\vec{F}_1+\vec{F}_2=\overrightarrow{MA}+\overrightarrow{MB}=\overrightarrow{MC}$ (với $C$ là điểm sao cho $AMBC$ là hình bình)\\
	Khi đó cường độ lực tác dụng lên vật là $\left|\vec{F}_1+\vec{F}_2\right|=|\overrightarrow{MC}|=MC$.\\
	Ta có $MA=\left|\overrightarrow{MA}\right|=\left|\vec{F}_1\right|=400$ N, $MB=\left|\overrightarrow{MB}\right|=\left|\vec{F}_2\right|=300$ N.\\
	Mặt khác, do $\widehat{AMB}=90^{\circ}$ nên $AMBC$ là hình chữ nhật.\\
	Khi đó $MC=\sqrt{MA^2+MB^2}=\sqrt{400^2+300^2}=500$ N.}
\end{ex}

\begin{ex}%[0H5V2-6]
	Cho ba lực $\vec{F}_1=\overrightarrow{MA}$, $\vec{F}_2=\overrightarrow{MB}$, $\vec{F}_3=\overrightarrow{MC}$ cùng tác động vào một ô tô tại điểm $M$ và ô tô đứng yên. Cho biết cường độ hai lực $\vec{F}_1$, $\vec{F}_2$ đều bằng $25$ N và góc $\widehat{AMB}=60^\circ$. Khi đó cường độ $\vec{F}_3$ bằng bao nhiêu Newton (làm tròn kết quả đến hàng phần chục của Newton)?
	\begin{center}
		\begin{tikzpicture}[line join=round,line cap=round,>=stealth,scale=1,font=\footnotesize]
			\foreach \x/\y/\n in {0/0/M} \coordinate (\n) at (\x,\y);
			\coordinate (A) at ($(M)+({2*cos(150)},{2*sin(150)})$);
			\coordinate (B) at ($(M)+({2*cos(210)},{2*sin(210)})$);
			\coordinate (D) at ($(A)+(B)-(M)$);
			\coordinate (C) at ($(M)+(M)-(D)$);
			\shade[ball color=white] (M) circle (.2cm);
			\draw[->,thick] (M)--(A) node[sloped,midway,above]{$\vec{F}_1$};
			\draw[->,thick] (M)--(B) node[sloped,midway,below]{$\vec{F}_2$};
			\draw[->,thick] (M)--(C) node[sloped,midway,above]{$\vec{F}_3$};
			\foreach \t/\g in {A/150,B/-150,M/-90,C/0}{
				\draw[fill=black] (\t) circle (1pt) node[shift={(\g:10pt)}]{$ \t $};
			}
			\draw pic[thin,draw,"$60^\circ$",angle radius=.8cm] {angle= A--M--B};
		\end{tikzpicture}
	\end{center}
	\shortans{$43{,}3$}
\loigiai{
	\begin{center}
		\begin{tikzpicture}[line join=round,line cap=round,>=stealth,scale=1,font=\footnotesize]
			\foreach \x/\y/\n in {0/0/M} \coordinate (\n) at (\x,\y);
			\coordinate (A) at ($(M)+({2*cos(150)},{2*sin(150)})$);
			\coordinate (B) at ($(M)+({2*cos(210)},{2*sin(210)})$);
			\coordinate (D) at ($(A)+(B)-(M)$);
			\coordinate (C) at ($(M)+(M)-(D)$);
			\shade[ball color=white] (M) circle (.2cm);
			\draw[->,thick] (M)--(A) node[sloped,midway,above]{$\vec{F}_1$};
			\draw[->,thick] (M)--(B) node[sloped,midway,below]{$\vec{F}_2$};
			\draw[->,thick] (M)--(C) node[sloped,midway,above]{$\vec{F}_3$};
			\draw[->,thick] (M)--(D);
			\draw[dashed] (A)--(D) (B)--(D);
			\foreach \t/\g in {A/150,B/-150,M/-90,C/0,D/180}{
				\draw[fill=black] (\t) circle (1pt) node[shift={(\g:10pt)}]{$ \t $};
			}
			\draw pic[thin,draw,"$60^\circ$",angle radius=.8cm] {angle= A--M--B};
		\end{tikzpicture}
	\end{center}
	Ta có $\vec{F}_1+\vec{F}_2=\overrightarrow{MA}+\overrightarrow{MB}=\overrightarrow{MD}$ (với $D$ là điểm sao cho $AMBD$ là hình bình hành).\\
	Ta có $MA=\left|\overrightarrow{MA}\right|=\left|\vec{F}_1\right|=25$ N và $MB=\left|\overrightarrow{MB}\right|=\left|\vec{F}_2\right|=25$ N.\\
	Do $\widehat{AMB}=60^{\circ}$ nên $\triangle MAB$ là tam giác đều. Khi đó $MD=2\cdot\dfrac{25\sqrt{3}}{2}=25\sqrt{3}$ N.\\
	Do ô tô đứng yên nên cường độ lực tác dụng lên ô tô bằng $0$ hay $\vec{F}_1+\vec{F}_2+\vec{F}_3=\vec{0}$.\\
	Suy ra $\vec{F}_3=-\left(\vec{F}_1+\vec{F}_2\right)\Rightarrow\left|\vec{F}_3\right|=\left|-\left(\vec{F}_1+\vec{F}_2\right)\right|=\left|\overrightarrow{DM}\right|=MD=25\sqrt{3}\approx 43{,}3$ N.\\
	Vậy cường độ của $\vec{F}_3\approx 43{,}3$ N.}
\end{ex}

\begin{ex}%[0H5V2-6]
	Cho ba lực $\vec{F}_1=\overrightarrow{MA}$, $\vec{F}_2=\overrightarrow{MB}$, $\vec{F}_3=\overrightarrow{MC}$ cùng tác động vào một vật tại điểm $M$ và vật đứng yên. Cho biết cường độ của $\vec{F}_1$, $\vec{F}_2$ đều bằng $100$ N và góc $\widehat{AMB}=90^{\circ}$. Khi đó cường độ $\vec{F}_3$ bằng bao nhiêu Newton (làm tròn kết quả đến hàng đơn vị của Newton)?
	\begin{center}
		\begin{tikzpicture}[line join=round,line cap=round,>=stealth,scale=1,font=\footnotesize]
			\foreach \x/\y/\n in {0/0/M} \coordinate (\n) at (\x,\y);
			\coordinate (A) at ($(M)+({2*cos(135)},{2*sin(135)})$);
			\coordinate (B) at ($(M)+({2*cos(225)},{2*sin(225)})$);
			\coordinate (D) at ($(A)+(B)-(M)$);
			\coordinate (C) at ($(M)+(M)-(D)$);
			\shade[ball color=white] (M) circle (.2cm);
			\draw[->,thick] (M)--(A) node[sloped,midway,above]{$\vec{F}_1$};
			\draw[->,thick] (M)--(B) node[sloped,midway,below]{$\vec{F}_2$};
			\draw[->,thick] (M)--(C) node[sloped,midway,above]{$\vec{F}_3$};
			\foreach \t/\g in {A/150,B/-150,M/-90,C/0}{
				\draw[fill=black] (\t) circle (1pt) node[shift={(\g:10pt)}]{$ \t $};
			}
			\draw pic[thin,draw,angle radius=3mm] {right angle= A--M--B};
		\end{tikzpicture}
	\end{center}
	\shortans{$141$}
\loigiai{
	Vẽ hình vuông $ADBM$, cạnh $MA=100$ thì $MD=100\sqrt{2}$. 
	\begin{center}
		\begin{tikzpicture}[line join=round,line cap=round,>=stealth,scale=1,font=\footnotesize]
			\foreach \x/\y/\n in {0/0/M} \coordinate (\n) at (\x,\y);
			\coordinate (A) at ($(M)+({2*cos(135)},{2*sin(135)})$);
			\coordinate (B) at ($(M)+({2*cos(225)},{2*sin(225)})$);
			\coordinate (D) at ($(A)+(B)-(M)$);
			\coordinate (C) at ($(M)+(M)-(D)$);
			\shade[ball color=white] (M) circle (.2cm);
			\draw[->,thick] (M)--(A) node[sloped,midway,above]{$\vec{F}_1$};
			\draw[->,thick] (M)--(B) node[sloped,midway,below]{$\vec{F}_2$};
			\draw[->,thick] (M)--(C) node[sloped,midway,above]{$\vec{F}_3$};
			\draw[->,thick] (M)--(D);
			\draw[dashed] (A)--(D) (B)--(D);
			\foreach \t/\g in {A/150,B/-150,M/-90,C/0,D/180}{
				\draw[fill=black] (\t) circle (1pt) node[shift={(\g:10pt)}]{$ \t $};
			}
			\draw pic[thin,draw,angle radius=3mm] {right angle= A--M--B};
		\end{tikzpicture}
	\end{center}
	Theo qui tắc hình bình hành có $\overrightarrow{MA}+\overrightarrow{MB}=\overrightarrow{MD}$. Vật đứng yên khi $\vec{F}_3=-\overrightarrow{MD}$.\\ 
	Suy ra cường độ của lực $\vec{F}_3$ là $\left|\vec{F}_3\right|=100\sqrt{2}\approx 141$ N.
	}
\end{ex}



\begin{ex}%[0H5H2-4]
	Cho hình thang $ABCD$ có hai đáy $AB=1$, $CD=2$. Gọi $M$, $N$ lần lượt là trung điểm $AD$ và $BC$. Tính $\left|\overrightarrow{DM}-\overrightarrow{BA}-\overrightarrow{CN}\right|$.
	\shortans{$1{,}5$}
\loigiai{
	Ta có
	$$\begin{aligned}
		\left|\overrightarrow{DM}-\overrightarrow{BA}-\overrightarrow{CN}\right|&=\left|\overrightarrow{MA}+\overrightarrow{AB}+\overrightarrow{NC}\right|= \left|\overrightarrow{MA}+\overrightarrow{AB}+\overrightarrow{BN}\right|\\
		&=\left|\overrightarrow{MN}\right|=MN=\dfrac{AB+CD}{2}=\dfrac{3a}{2}.
	\end{aligned}$$
	}
\end{ex}

\begin{ex}%[0H5H2-4]
	Cho tam giác $ABC$ vuông tại $A$ có $\widehat{ABC}=30^{\circ}$ và $BC=\sqrt{5}$.
	Tính độ dài của véc-tơ $\overrightarrow{AB}+\overrightarrow{AC}$ (làm tròn kết quả đến hàng phần trăm). 
	\shortans{$2{,}24$}
\loigiai{
	\immini{Gọi $D$ là điểm sao cho tứ giác $ABDC$ là hình bình hành.\\
	Khi đó theo quy tắc hình bình hành ta có $\overrightarrow{AB}+\overrightarrow{AC}=\overrightarrow{AD}$.\\
	Vì $\triangle ABC$ vuông ở $A$ nên tứ giác $ABDC$ là hình chữ nhật.\\
	Suy ra $AD=BC=a\sqrt{5}$.\\
	Vậy $\left|\overrightarrow{AB}+\overrightarrow{AC}\right|=\left|\overrightarrow{AD}\right|=AD=a\sqrt{5}$.
	}{\begin{tikzpicture}[line join=round,line cap=round,>=stealth,scale=1,font=\footnotesize]
			\foreach \x/\y/\n in {0/0/A,3/0/B,0/1.5/C} \coordinate (\n) at (\x,\y);
			\coordinate (D) at ($(B)+(C)-(A)$); 
			\draw (A)--(B)--(C)--(D)--(A)--(C) (B)--(D);
			\foreach \t/\g in {A/-135,D/45,C/135,B/-45}{
				\draw[fill=black] (\t) circle (1pt) node[shift={(\g:8pt)}]{$ \t $};
			}
		\end{tikzpicture}}
	}
\end{ex}







\Closesolutionfile{ans}
\indapan{6}{ans/ans-0-C4B2-2-KQ}